\begin{recipe}
[% 
    preparationtime = {\unit[15]{min} + \unit[2]{h} Kühlzeit},
    portion = {\portion{6}},
    source = {\protect\recipesource{https://veggie-einhorn.de/vegane-mousse-au-chocolat/}{Veggie Einhorn: Vegane Mousse au Chocolat mit Seidentofu}},
    calory = {349kcal | 5g Protein | 31g KH | 21g Fett},
    price = {1,05€},
    created = {05.05.2026},
    rating = {4},
]
{Vegane Mousse au Chocolat}

\graph{% pictures
    big=pic/dessert/10_vegane-mousse-au-chocolat
}

\introduction{%
Eine simple vegane Mousse au Chocolat aus Seidentofu und dunkler Schokolade. Sehr wenig Aufwand, aber trotzdem cremig, schokoladig und perfekt, wenn man Gäste beeindrucken will ohne wirklich viel zu machen.
}

\ingredients{
    400 g & Seidentofu\index{Tofu \& Fleischalternativen!Seidentofu}\\
    300 g & Vegane Zartbitterschokolade\index{Backzutaten \& Süßes!Schokolade}\\
    3 Päckchen & Vanillezucker\index{Backzutaten \& Süßes!Vanillezucker}
}

\preparation{%
    \step%
    Schokolade im Wasserbad langsam schmelzen. Dabei darauf achten, dass kein Wasser in die Schokolade kommt.
    
    \step%
    Seidentofu, Vanillezucker und die geschmolzene Schokolade in einen Mixer geben.
    
    \step%
    Alles gründlich mixen, bis eine glatte, cremige Masse entsteht.
    
    \step%
    Mousse in Dessertgläser oder eine Schüssel füllen und mindestens 2 Stunden kalt stellen.
}

\hint{
    Die Schokolade nicht zu heiß werden lassen, sonst kann sie grisselig werden. Lieber langsam schmelzen und kurz abkühlen lassen, bevor sie zum Seidentofu kommt.
}

\suggestion[Servieren]{%
    Mit frischen Beeren, gehackter Schokolade oder etwas Kakaopulver servieren.
}

\end{recipe}